% =============================================================================
% Section: 02-problem-and-use-case.tex | Rubric: Problema y caso de uso
% Purpose: Problema correcto, costo de no actuar y caso de uso acotado.
% Style: es-CO académico-profesional, sin sobrepromesas; clasificar claims como
%        implementado/demo/pendiente; labels fig:/tab: estables (snake_case).
% =============================================================================
\section{Problema y caso de uso}

El problema correcto no es que las PyMEs carezcan de un canal digital. El canal ya existe: WhatsApp, cuya plataforma de negocio ofrece integración programática mediante la Cloud API de Meta~\cite{meta2026whatsapp}. El problema es que el proceso comercial ocurre sobre mensajes no estructurados. En pedidos personalizados, la operación necesita datos mínimos: producto, cantidad, fecha, dirección, contacto, observaciones y, en algunos casos, una imagen de referencia. Cuando esos datos llegan incompletos, el negocio queda atrapado en ciclos de aclaración, reproceso y decisiones informales.

El proceso actual puede resumirse en cuatro pasos. Primero, el cliente escribe como le resulta natural. Segundo, una persona interpreta el pedido y pregunta lo que falta. Tercero, otra persona u otra área produce con la información disponible. Cuarto, el propietario o encargado absorbe errores, demoras y reclamos. El costo de no actuar no se limita a horas humanas: también incluye ventas perdidas por respuesta lenta, pedidos incompletos, baja trazabilidad y dificultad para medir cuántas conversaciones se convierten realmente en pedidos.

El caso de uso de DeliverAI se centra en una pastelería o negocio similar con pedidos personalizados. Este dominio es apropiado para un piloto porque combina lenguaje natural, imágenes, aclaraciones y datos transaccionales, pero permite acotar el catálogo, los horarios y las reglas. Además, el riesgo puede controlarse: el agente no cobra, no promete entregas fuera de política y no registra pedidos sin confirmación explícita del cliente.
